% Deck 0 of the workshop: setting up before the hands-on sessions.
% Installation of Git, Ollama, OpenCode, Claude Code and Codex; Git in ten
% commands; the controls of a harness. Commands were taken from the tools'
% own documentation and from the installed tools on 2 November 2026; the
% same text, for copying, is in workshop/SETUP.md.
% Build: latexmk -pdf main.tex   (pdflatex + biber)
\documentclass[11pt,aspectratio=169]{beamer}
% ---------------------------------------------------------------------------
% preamble.tex --- shared Beamer style for the four workshop decks
%
% Loaded by each deck with % ---------------------------------------------------------------------------
% preamble.tex --- shared Beamer style for the four workshop decks
%
% Loaded by each deck with % ---------------------------------------------------------------------------
% preamble.tex --- shared Beamer style for the four workshop decks
%
% Loaded by each deck with \input{../common/preamble}. Same fonts, colours
% and bibliography settings as the book (book/preamble.tex): Latin Modern,
% black on white, one grey for code, one accent colour for what is untrusted
% or wrong. Built with pdflatex + biber, like the book.
% ---------------------------------------------------------------------------

\usepackage[T1]{fontenc}
\usepackage[utf8]{inputenc}
\usepackage{lmodern}
\usepackage{microtype}
\usepackage{booktabs}
\usepackage{tabularx}
\usepackage{listings}
\usepackage{xcolor}
\usepackage{csquotes}
\usepackage{graphicx}

% Figures: the book's figures as the 300 dpi PNGs the book itself includes
% (book/figures/png, tracked in git and current; figures/pdf is not rebuilt
% after every change), and the first pages of the papers
% (common/screenshots). Both are found by name without an extension.
\graphicspath{{../../book/figures/png/}{../common/screenshots/}}
\DeclareGraphicsExtensions{.png,.pdf}

% --- look ------------------------------------------------------------------
\usetheme{default}
\usefonttheme{serif}
\definecolor{accent}{HTML}{9B2D20}
\definecolor{codebg}{gray}{0.96}
\definecolor{rule}{gray}{0.75}
\definecolor{dim}{gray}{0.45}
\setbeamercolor{structure}{fg=black}
\setbeamercolor{normal text}{fg=black,bg=white}
\setbeamercolor{alerted text}{fg=accent}
\setbeamercolor{section in toc}{fg=black}
\setbeamercolor{footline}{fg=dim}
\setbeamerfont{title}{series=\bfseries,size=\Large}
\setbeamerfont{subtitle}{series=\mdseries,size=\normalsize}
\setbeamerfont{frametitle}{series=\bfseries,size=\large}
\setbeamerfont{footline}{size=\scriptsize}
\setbeamertemplate{navigation symbols}{}
\setbeamertemplate{itemize items}[circle]
\setbeamertemplate{enumerate items}[default]
\setbeamertemplate{caption}{\scriptsize\insertcaption}
\setbeamertemplate{frametitle}{%
  \vskip0.6em\usebeamerfont{frametitle}\insertframetitle\par
  \vskip0.2em{\color{rule}\hrule height 0.4pt}\vskip0.4em}
\setbeamertemplate{footline}{%
  \hbox to \paperwidth{\hspace{1.2em}\usebeamerfont{footline}\usebeamercolor[fg]{footline}%
    \insertshorttitle\ --- \insertsection\hfill
    \insertframenumber/\inserttotalframenumber\hspace{1.2em}}\vskip0.6em}
\setbeamersize{text margin left=1.4em,text margin right=1.4em}
% An outline frame at the start of each section, the current one in black.
\AtBeginSection[]{\begin{frame}{Outline}\tableofcontents[currentsection]\end{frame}}
% A frame that breaks over pages is numbered only from its second page.
\setbeamertemplate{frametitle continuation}[from second]

% --- code listings (as in the book) ----------------------------------------
\lstset{
  basicstyle=\ttfamily\footnotesize,
  backgroundcolor=\color{codebg},
  frame=none,
  breaklines=true,
  breakatwhitespace=true,
  columns=fullflexible,
  keepspaces=true,
  upquote=true,
  showstringspaces=false,
  xleftmargin=0.6em,
  framexleftmargin=0.6em,
  aboveskip=0.5em,
  belowskip=0.5em,
  literate={~}{{\textasciitilde}}1
           {>}{{\textgreater}}1
           {<}{{\textless}}1
           {ö}{{\"o}}1 {ü}{{\"u}}1 {é}{{\'e}}1,
}
\lstdefinestyle{shell}{language=bash}

% --- bibliography (same settings as the book) ------------------------------
\usepackage[
  backend=biber,
  style=authoryear,
  sorting=nyt,
  maxbibnames=6,
  minbibnames=3,
  maxcitenames=2,
  giveninits=true,
  uniquename=init,
  doi=true,
  url=true,
  isbn=false,
  eprint=true,
]{biblatex}
\addbibresource{../../book/references.bib}
\renewcommand*{\bibfont}{\scriptsize}
\setlength{\bibitemsep}{0.25em}
\setbeamertemplate{bibliography item}{}   % no icon before each entry

% --- shorthand -------------------------------------------------------------
\newcommand{\code}[1]{\texttt{#1}}
\newcommand{\term}[1]{\emph{#1}}
% A key point, as in the book: no box.
\newcommand{\kp}[1]{\par\medskip{\bfseries Key point.}\enspace\emph{#1}\par}
% \fig[height]{name}: one of the book's figures, scaled to fit the frame.
\newcommand{\fig}[2][0.72\textheight]{%
  \centering\includegraphics[height=#1,width=\textwidth,keepaspectratio]{#2}\par}
% \screenshot{key}: the first page of a paper, with a thin frame.
\newcommand{\screenshot}[2][0.66\textheight]{%
  \setlength{\fboxsep}{0pt}\setlength{\fboxrule}{0.3pt}%
  \fbox{\includegraphics[height=#1,width=\linewidth,keepaspectratio]{#2}}}
% \paperslide{key}{text}: screenshot on the left, what the book takes from the
% paper on the right. Use inside a frame; add \source after it.
\newcommand{\paperslide}[2]{%
  \begin{columns}[T]
    \begin{column}{0.40\textwidth}\centering\screenshot{#1}\end{column}
    \begin{column}{0.58\textwidth}\small #2\end{column}
  \end{columns}}
% Notes in references.bib record how a web source was saved and dated; they
% belong in the reference list, not on a \source line, so cite commands
% (including \fullcite) leave them out.
\AtEveryCitekey{\clearfield{note}}
% \source{...}: a small grey line at the foot of a frame, for the source.
\newcommand{\source}[1]{\vskip0pt plus 1filll\par{\scriptsize\color{dim}#1\par}}
. Same fonts, colours
% and bibliography settings as the book (book/preamble.tex): Latin Modern,
% black on white, one grey for code, one accent colour for what is untrusted
% or wrong. Built with pdflatex + biber, like the book.
% ---------------------------------------------------------------------------

\usepackage[T1]{fontenc}
\usepackage[utf8]{inputenc}
\usepackage{lmodern}
\usepackage{microtype}
\usepackage{booktabs}
\usepackage{tabularx}
\usepackage{listings}
\usepackage{xcolor}
\usepackage{csquotes}
\usepackage{graphicx}

% Figures: the book's figures as the 300 dpi PNGs the book itself includes
% (book/figures/png, tracked in git and current; figures/pdf is not rebuilt
% after every change), and the first pages of the papers
% (common/screenshots). Both are found by name without an extension.
\graphicspath{{../../book/figures/png/}{../common/screenshots/}}
\DeclareGraphicsExtensions{.png,.pdf}

% --- look ------------------------------------------------------------------
\usetheme{default}
\usefonttheme{serif}
\definecolor{accent}{HTML}{9B2D20}
\definecolor{codebg}{gray}{0.96}
\definecolor{rule}{gray}{0.75}
\definecolor{dim}{gray}{0.45}
\setbeamercolor{structure}{fg=black}
\setbeamercolor{normal text}{fg=black,bg=white}
\setbeamercolor{alerted text}{fg=accent}
\setbeamercolor{section in toc}{fg=black}
\setbeamercolor{footline}{fg=dim}
\setbeamerfont{title}{series=\bfseries,size=\Large}
\setbeamerfont{subtitle}{series=\mdseries,size=\normalsize}
\setbeamerfont{frametitle}{series=\bfseries,size=\large}
\setbeamerfont{footline}{size=\scriptsize}
\setbeamertemplate{navigation symbols}{}
\setbeamertemplate{itemize items}[circle]
\setbeamertemplate{enumerate items}[default]
\setbeamertemplate{caption}{\scriptsize\insertcaption}
\setbeamertemplate{frametitle}{%
  \vskip0.6em\usebeamerfont{frametitle}\insertframetitle\par
  \vskip0.2em{\color{rule}\hrule height 0.4pt}\vskip0.4em}
\setbeamertemplate{footline}{%
  \hbox to \paperwidth{\hspace{1.2em}\usebeamerfont{footline}\usebeamercolor[fg]{footline}%
    \insertshorttitle\ --- \insertsection\hfill
    \insertframenumber/\inserttotalframenumber\hspace{1.2em}}\vskip0.6em}
\setbeamersize{text margin left=1.4em,text margin right=1.4em}
% An outline frame at the start of each section, the current one in black.
\AtBeginSection[]{\begin{frame}{Outline}\tableofcontents[currentsection]\end{frame}}
% A frame that breaks over pages is numbered only from its second page.
\setbeamertemplate{frametitle continuation}[from second]

% --- code listings (as in the book) ----------------------------------------
\lstset{
  basicstyle=\ttfamily\footnotesize,
  backgroundcolor=\color{codebg},
  frame=none,
  breaklines=true,
  breakatwhitespace=true,
  columns=fullflexible,
  keepspaces=true,
  upquote=true,
  showstringspaces=false,
  xleftmargin=0.6em,
  framexleftmargin=0.6em,
  aboveskip=0.5em,
  belowskip=0.5em,
  literate={~}{{\textasciitilde}}1
           {>}{{\textgreater}}1
           {<}{{\textless}}1
           {ö}{{\"o}}1 {ü}{{\"u}}1 {é}{{\'e}}1,
}
\lstdefinestyle{shell}{language=bash}

% --- bibliography (same settings as the book) ------------------------------
\usepackage[
  backend=biber,
  style=authoryear,
  sorting=nyt,
  maxbibnames=6,
  minbibnames=3,
  maxcitenames=2,
  giveninits=true,
  uniquename=init,
  doi=true,
  url=true,
  isbn=false,
  eprint=true,
]{biblatex}
\addbibresource{../../book/references.bib}
\renewcommand*{\bibfont}{\scriptsize}
\setlength{\bibitemsep}{0.25em}
\setbeamertemplate{bibliography item}{}   % no icon before each entry

% --- shorthand -------------------------------------------------------------
\newcommand{\code}[1]{\texttt{#1}}
\newcommand{\term}[1]{\emph{#1}}
% A key point, as in the book: no box.
\newcommand{\kp}[1]{\par\medskip{\bfseries Key point.}\enspace\emph{#1}\par}
% \fig[height]{name}: one of the book's figures, scaled to fit the frame.
\newcommand{\fig}[2][0.72\textheight]{%
  \centering\includegraphics[height=#1,width=\textwidth,keepaspectratio]{#2}\par}
% \screenshot{key}: the first page of a paper, with a thin frame.
\newcommand{\screenshot}[2][0.66\textheight]{%
  \setlength{\fboxsep}{0pt}\setlength{\fboxrule}{0.3pt}%
  \fbox{\includegraphics[height=#1,width=\linewidth,keepaspectratio]{#2}}}
% \paperslide{key}{text}: screenshot on the left, what the book takes from the
% paper on the right. Use inside a frame; add \source after it.
\newcommand{\paperslide}[2]{%
  \begin{columns}[T]
    \begin{column}{0.40\textwidth}\centering\screenshot{#1}\end{column}
    \begin{column}{0.58\textwidth}\small #2\end{column}
  \end{columns}}
% Notes in references.bib record how a web source was saved and dated; they
% belong in the reference list, not on a \source line, so cite commands
% (including \fullcite) leave them out.
\AtEveryCitekey{\clearfield{note}}
% \source{...}: a small grey line at the foot of a frame, for the source.
\newcommand{\source}[1]{\vskip0pt plus 1filll\par{\scriptsize\color{dim}#1\par}}
. Same fonts, colours
% and bibliography settings as the book (book/preamble.tex): Latin Modern,
% black on white, one grey for code, one accent colour for what is untrusted
% or wrong. Built with pdflatex + biber, like the book.
% ---------------------------------------------------------------------------

\usepackage[T1]{fontenc}
\usepackage[utf8]{inputenc}
\usepackage{lmodern}
\usepackage{microtype}
\usepackage{booktabs}
\usepackage{tabularx}
\usepackage{listings}
\usepackage{xcolor}
\usepackage{csquotes}
\usepackage{graphicx}

% Figures: the book's figures as the 300 dpi PNGs the book itself includes
% (book/figures/png, tracked in git and current; figures/pdf is not rebuilt
% after every change), and the first pages of the papers
% (common/screenshots). Both are found by name without an extension.
\graphicspath{{../../book/figures/png/}{../common/screenshots/}}
\DeclareGraphicsExtensions{.png,.pdf}

% --- look ------------------------------------------------------------------
\usetheme{default}
\usefonttheme{serif}
\definecolor{accent}{HTML}{9B2D20}
\definecolor{codebg}{gray}{0.96}
\definecolor{rule}{gray}{0.75}
\definecolor{dim}{gray}{0.45}
\setbeamercolor{structure}{fg=black}
\setbeamercolor{normal text}{fg=black,bg=white}
\setbeamercolor{alerted text}{fg=accent}
\setbeamercolor{section in toc}{fg=black}
\setbeamercolor{footline}{fg=dim}
\setbeamerfont{title}{series=\bfseries,size=\Large}
\setbeamerfont{subtitle}{series=\mdseries,size=\normalsize}
\setbeamerfont{frametitle}{series=\bfseries,size=\large}
\setbeamerfont{footline}{size=\scriptsize}
\setbeamertemplate{navigation symbols}{}
\setbeamertemplate{itemize items}[circle]
\setbeamertemplate{enumerate items}[default]
\setbeamertemplate{caption}{\scriptsize\insertcaption}
\setbeamertemplate{frametitle}{%
  \vskip0.6em\usebeamerfont{frametitle}\insertframetitle\par
  \vskip0.2em{\color{rule}\hrule height 0.4pt}\vskip0.4em}
\setbeamertemplate{footline}{%
  \hbox to \paperwidth{\hspace{1.2em}\usebeamerfont{footline}\usebeamercolor[fg]{footline}%
    \insertshorttitle\ --- \insertsection\hfill
    \insertframenumber/\inserttotalframenumber\hspace{1.2em}}\vskip0.6em}
\setbeamersize{text margin left=1.4em,text margin right=1.4em}
% An outline frame at the start of each section, the current one in black.
\AtBeginSection[]{\begin{frame}{Outline}\tableofcontents[currentsection]\end{frame}}
% A frame that breaks over pages is numbered only from its second page.
\setbeamertemplate{frametitle continuation}[from second]

% --- code listings (as in the book) ----------------------------------------
\lstset{
  basicstyle=\ttfamily\footnotesize,
  backgroundcolor=\color{codebg},
  frame=none,
  breaklines=true,
  breakatwhitespace=true,
  columns=fullflexible,
  keepspaces=true,
  upquote=true,
  showstringspaces=false,
  xleftmargin=0.6em,
  framexleftmargin=0.6em,
  aboveskip=0.5em,
  belowskip=0.5em,
  literate={~}{{\textasciitilde}}1
           {>}{{\textgreater}}1
           {<}{{\textless}}1
           {ö}{{\"o}}1 {ü}{{\"u}}1 {é}{{\'e}}1,
}
\lstdefinestyle{shell}{language=bash}

% --- bibliography (same settings as the book) ------------------------------
\usepackage[
  backend=biber,
  style=authoryear,
  sorting=nyt,
  maxbibnames=6,
  minbibnames=3,
  maxcitenames=2,
  giveninits=true,
  uniquename=init,
  doi=true,
  url=true,
  isbn=false,
  eprint=true,
]{biblatex}
\addbibresource{../../book/references.bib}
\renewcommand*{\bibfont}{\scriptsize}
\setlength{\bibitemsep}{0.25em}
\setbeamertemplate{bibliography item}{}   % no icon before each entry

% --- shorthand -------------------------------------------------------------
\newcommand{\code}[1]{\texttt{#1}}
\newcommand{\term}[1]{\emph{#1}}
% A key point, as in the book: no box.
\newcommand{\kp}[1]{\par\medskip{\bfseries Key point.}\enspace\emph{#1}\par}
% \fig[height]{name}: one of the book's figures, scaled to fit the frame.
\newcommand{\fig}[2][0.72\textheight]{%
  \centering\includegraphics[height=#1,width=\textwidth,keepaspectratio]{#2}\par}
% \screenshot{key}: the first page of a paper, with a thin frame.
\newcommand{\screenshot}[2][0.66\textheight]{%
  \setlength{\fboxsep}{0pt}\setlength{\fboxrule}{0.3pt}%
  \fbox{\includegraphics[height=#1,width=\linewidth,keepaspectratio]{#2}}}
% \paperslide{key}{text}: screenshot on the left, what the book takes from the
% paper on the right. Use inside a frame; add \source after it.
\newcommand{\paperslide}[2]{%
  \begin{columns}[T]
    \begin{column}{0.40\textwidth}\centering\screenshot{#1}\end{column}
    \begin{column}{0.58\textwidth}\small #2\end{column}
  \end{columns}}
% Notes in references.bib record how a web source was saved and dated; they
% belong in the reference list, not on a \source line, so cite commands
% (including \fullcite) leave them out.
\AtEveryCitekey{\clearfield{note}}
% \source{...}: a small grey line at the foot of a frame, for the source.
\newcommand{\source}[1]{\vskip0pt plus 1filll\par{\scriptsize\color{dim}#1\par}}


\title[AI in Practice 0]{AI in Practice, or: How I Learned to Stop Worrying and Love the Agent}
\subtitle{A toolbox and worked examples for biodiversity researchers, from the book \emph{AI for Biodiversity: A Primer}\\[0.6em] 0. Setting up: the tools, Git, and driving a harness}
\author{Savvas Paragkamian}
\institute{Hellenic Centre for Marine Research \\ KYBELE and ELIXIR-GR}
\date{HCMR workshop, November 2026}
\hypersetup{pdftitle={AI for Biodiversity: A Primer, deck 0},pdfauthor={Savvas Paragkamian}}

\begin{document}

\begin{frame}[plain]
  \titlepage
\end{frame}

\begin{frame}{Goals of this session}
  \begin{enumerate}
    \item Install what the hands-on sessions need and check that each piece
      answers: Git, Ollama with one model, and at least one harness.
    \item Learn the ten Git commands that make an agent session safe: a
      checkpoint before, a diff after, and a way back.
    \item Learn the controls every harness has: which model, how much
      reasoning, how full the context is, what it may touch, and how to run
      it without a conversation.
  \end{enumerate}
  \kp{Commands and versions on these slides were checked on 2 November 2026
  and will date; the reasoning behind them should not. The same commands,
  for copying, are in \code{workshop/SETUP.md}.}
\end{frame}

\section{What you need}

\begin{frame}{Three layers, and what sits in each}
  \begin{itemize}
    \item \textbf{The operating system and a terminal.} macOS and Linux as
      they are. On Windows, the book's advice (Chapter 8) is a Linux
      environment inside Windows (WSL); the harnesses also run natively.
    \item \textbf{R}, for the exercises: base R is enough for the agent
      session; the testing exercise needs the \code{ellmer} package.
    \item \textbf{Git}, for the checkpoint habit.
    \item \textbf{A runtime for local models:} Ollama, and one model pulled
      into it (the local track).
    \item \textbf{A harness:} OpenCode (open source, works with Ollama),
      Claude Code (hosted; needs a paid Anthropic plan) or Codex (needs a
      ChatGPT plan, or runs local models through Ollama).
  \end{itemize}
  \source{Chapter 8, The Software Stack; workshop/README.md, ``What you need''.}
\end{frame}

\begin{frame}{The versions this workshop was prepared with}
  \begin{tabularx}{\textwidth}{lXl}
    \toprule
    \textbf{Tool} & \textbf{What it is} & \textbf{Version, 2 Nov 2026} \\
    \midrule
    Git & version control & 2.54 (Apple); 2.56.0 at git-scm.com \\
    Ollama & runs open-weight models locally & 0.40.13 \\
    OpenCode & open-source terminal harness (MIT) & 1.17.11 \\
    Claude Code & Anthropic's terminal harness & 2.1.316 \\
    Codex & OpenAI's terminal harness (Apache-2.0) & 0.162.1 \\
    \bottomrule
  \end{tabularx}
  \medskip
  Names and versions date within months; Appendix D of the book keeps the
  dated details, and each tool prints its own with \code{--version}.
  \source{Versions read from the installed tools on the author's machine; licences from Appendix D.}
\end{frame}

\section{Installing}

\begin{frame}[fragile]{Git}
  \begin{lstlisting}[style=shell]
# macOS: Apple's command line tools, or Homebrew
xcode-select --install
brew install git
# Debian and Ubuntu                    # Fedora
sudo apt-get install git               sudo dnf install git
# Windows: the Git for Windows installer from git-scm.com, or
winget install --id Git.Git -e --source winget
# then, in any terminal
git --version
git config --global user.name  "Your Name"
git config --global user.email "you@example.org"
  \end{lstlisting}
  Git for Windows also gives the harnesses a Bash shell to run commands in.
  \source{git-scm.com/install (read 2 November 2026); Claude Code setup notes on Git for Windows.}
\end{frame}

\begin{frame}[fragile]{Ollama, and a model}
  \begin{lstlisting}[style=shell]
# macOS and Linux (or the .dmg from ollama.com/download)
curl -fsSL https://ollama.com/install.sh | sh
# Windows (PowerShell), or OllamaSetup.exe from ollama.com/download
irm https://ollama.com/install.ps1 | iex
ollama --version
ollama pull gemma4:12b      # the model workshop/session/opencode.json names
ollama list                 # what is on disk
ollama ps                   # what is loaded, and with which CONTEXT
  \end{lstlisting}
  \begin{itemize}
    \item A 12B model at four bits needs about 8\,GB of memory; the context
      adds more (Chapter 13). A laptop with 16\,GB runs it; 32\,GB runs it
      comfortably.
    \item \code{ollama run MODEL} starts a chat; \code{/bye} leaves it.
  \end{itemize}
  \source{ollama.com and the Ollama README (read 2 November 2026); workshop/README.md for the model tags used in the test runs.}
\end{frame}

\begin{frame}[fragile]{Ollama: the context, the setting that breaks agents}
  \begin{lstlisting}[style=shell]
# once, for the whole server (or the slider in the Ollama app's settings)
OLLAMA_CONTEXT_LENGTH=32768 ollama serve
# or per model, saved under a new name
ollama run MODEL
  /set parameter num_ctx 32768
  /save MODEL-32k
ollama ps                   # CONTEXT must show the new value
  \end{lstlisting}
  \begin{itemize}
    \item Ollama's default context depends on the graphics memory: 4K below
      24\,GiB. An agent re-sends the whole session every turn and needs tens
      of thousands of tokens; at 4K it forgets the task and looks broken.
    \item A larger context needs more memory. Start at 32K; the vendor
      advises 64K for agents.
  \end{itemize}
  \source{Chapter 8; Appendix D; docs.ollama.com/context-length (read 2 November 2026).}
\end{frame}

\begin{frame}[fragile]{OpenCode}
  \begin{lstlisting}[style=shell]
curl -fsSL https://opencode.ai/install | bash     # macOS, Linux, WSL
npm install -g opencode-ai                         # any system with Node
brew install anomalyco/tap/opencode                # Homebrew, version 1
scoop install opencode                             # Windows (or choco)
opencode --version
cd ~/survey && opencode                            # start in the project
  \end{lstlisting}
  \begin{itemize}
    \item Hosted providers: type \code{/connect} in the TUI and paste a key;
      it is stored under \code{\textasciitilde/.local/share/opencode/}.
    \item Local models: a \code{provider} block in the project's
      \code{opencode.json} points at Ollama (next slide). No account needed.
    \item \code{opencode upgrade} updates it.
  \end{itemize}
  \source{opencode.ai/docs (read 2 November 2026). The workshop was tested with version 1.17.11.}
\end{frame}

\begin{frame}[fragile]{OpenCode and Ollama: the project's \code{opencode.json}}
  \begin{lstlisting}
{ "$schema": "https://opencode.ai/config.json",
  "model": "ollama/gemma4:12b",
  "provider": { "ollama": {
      "npm": "@ai-sdk/openai-compatible",
      "name": "Ollama (local)",
      "options": { "baseURL": "http://localhost:11434/v1" },
      "models": { "gemma4:12b": {
          "name": "Gemma 4 12B (local)",
          "tool_call": true,
          "limit": { "context": 32768, "output": 8192 } } } } },
  "permission": { ... } }         # the rules, shown in deck 3
  \end{lstlisting}
  Two settings cause most local-agent failures: the harness must be told the
  model supports tools (\code{tool\_call}), and the context must be enlarged
  on both sides (Chapter 8).
  \source{workshop/session/opencode.json, abridged; opencode.ai/docs/providers.}
\end{frame}

\begin{frame}[fragile]{Claude Code}
  \begin{lstlisting}[style=shell]
curl -fsSL https://claude.ai/install.sh | bash     # macOS, Linux, WSL
brew install --cask claude-code                    # Homebrew
irm https://claude.ai/install.ps1 | iex            # Windows PowerShell
winget install Anthropic.ClaudeCode                # Windows
claude --version
claude doctor                   # checks the installation and settings
cd ~/survey && claude           # first run opens the browser to log in
  \end{lstlisting}
  \begin{itemize}
    \item Needs a Pro, Max, Team or Enterprise plan, or API access; the free
      plan does not include it. Data goes to Anthropic: this is the hosted
      track.
    \item On Windows, the sandbox works only inside WSL~2.
    \item Native installs update themselves; \code{claude update} forces it.
  \end{itemize}
  \source{code.claude.com/docs/en/setup (read 2 November 2026).}
\end{frame}

\begin{frame}[fragile]{Codex}
  \begin{lstlisting}[style=shell]
npm install -g @openai/codex                       # any system with Node
brew install --cask codex                          # Homebrew
curl -fsSL https://chatgpt.com/codex/install.sh | sh   # macOS, Linux
codex --version
cd ~/survey && codex            # first run: "Sign in with ChatGPT"
codex --oss --local-provider ollama    # a local model through Ollama
  \end{lstlisting}
  \begin{itemize}
    \item Sign in with a ChatGPT plan, or with an API key
      (\code{codex login --with-api-key}).
    \item The settings live in \code{\textasciitilde/.codex/config.toml};
      \code{codex doctor} checks the installation.
    \item Codex was not used in the book's test runs; its commands here were
      read from its help and documentation.
  \end{itemize}
  \source{github.com/openai/codex and developers.openai.com/codex (read 2 November 2026); codex --help, version 0.162.1.}
\end{frame}

\section{Git in ten commands}

\begin{frame}{Why Git, before any agent starts}
  \begin{itemize}
    \item Use Git for code, always; with agents it stops being optional
      (Chapter 6).
    \item \textbf{Commit before} an agent starts: the commit is the checkpoint.
      Whatever happens next can be seen with \code{git diff} and undone with
      \code{git restore} (Chapter 10).
    \item \textbf{Read the diff after.} An agent's summary is a claim about
      its work; the diff is the evidence.
    \item Git protects only what is tracked and committed. Files an agent
      creates or deletes before the next commit are not recovered.
    \item An agent with a shell can itself run \code{git reset --hard},
      \code{git clean} or \code{git push --force}. Deny those in the
      harness configuration, as the workshop's settings do.
  \end{itemize}
  \source{Chapter 6, Version control; Chapter 10, Step 1.}
\end{frame}

\begin{frame}[fragile]{The ten commands}
  \begin{lstlisting}[style=shell]
git init                          # start tracking this folder
git status                        # what changed, what is untracked
git add -A                        # stage everything
git commit -m "Before the agent session"   # the checkpoint
git diff                          # what the agent changed, line by line
git log --oneline                 # the history, one line per commit
git restore scripts/01_check.R    # undo one file to the last commit
git restore .                     # undo every tracked file
git stash                         # park uncommitted changes, get them back with
                                  #   git stash pop
git remote add origin <url> && git push -u origin main   # a copy elsewhere
  \end{lstlisting}
  A \code{.gitignore} file lists what Git should not track: large data,
  build products, secrets. Never commit a key or a token.
  \source{Pro Git, chapter 2 (git-scm.com/book, free); Chapter 15 on credentials.}
\end{frame}

\begin{frame}[fragile]{Reading a diff after a session}
  \begin{lstlisting}[style=shell]
git status --short                # only scripts/, data/processed/, logs/ new?
git status --short data/raw       # must print nothing
git diff --stat                   # which files, how many lines
git diff data/                    # the lines themselves
git restore .                     # back to the checkpoint if it is beyond repair
  \end{lstlisting}
  \begin{itemize}
    \item Look first at what should not have changed: the raw data, the rules
      file, the configuration.
    \item Then at what the agent added: a script you can read and rerun, or
      an edit in place that you cannot?
    \item Commit the parts you accept, with a message that says what the
      agent did and what you checked.
  \end{itemize}
  \source{Chapter 10, Steps 4 to 6; Chapter 7, the six checks.}
\end{frame}

\section{Driving a harness}

\begin{frame}{The same five controls, in every harness}
  \small
  \begin{tabularx}{\textwidth}{lXXX}
    \toprule
    & \textbf{Claude Code} & \textbf{OpenCode} & \textbf{Codex} \\
    \midrule
    Model & \code{/model} & \code{/models} & \code{/model} \\
    Reasoning effort & \code{/effort}, or arrows in \code{/model} & \code{Ctrl+T} cycles the variant & in \code{/model} \\
    Context and usage & \code{/context}, \code{/usage} & \code{Ctrl+X S}; \code{opencode stats} & \code{/status} \\
    Permissions & \code{/permissions}; \code{Shift+Tab} cycles modes & \code{opencode.json}; \code{Tab} Plan/Build & \code{/permissions}; \code{-s}, \code{-a} \\
    Non-interactive & \code{claude -p "..."} & \code{opencode run "..."} & \code{codex exec "..."} \\
    Stop, compact & \code{Esc}; \code{/compact} & \code{Esc}; \code{/compact} & \code{/compact} \\
    Rules and memory & \code{CLAUDE.md} or \code{AGENTS.md}; \code{/memory} & \code{AGENTS.md}; \code{/init} & \code{AGENTS.md}; \code{/init} \\
    \bottomrule
  \end{tabularx}
  \medskip
  Type \code{/} in any of them to see the list; names change between
  versions, the controls do not.
  \source{The tools' documentation and \code{--help}, read 2 November 2026;
    the rules and memory row on 4 November 2026.}
\end{frame}

\begin{frame}{Claude Code: the controls}
  \begin{itemize}
    \item \code{/model} switches the model and saves it as the default;
      left and right arrows adjust the effort. \code{/effort low} to
      \code{xhigh}, \code{max} or \code{auto} sets it directly;
      \code{/fast} toggles fast mode.
    \item \code{/context} shows what fills the window; \code{/usage}
      (alias \code{/cost}) shows session cost, plan limits and activity;
      \code{/compact} summarises to free space; \code{/clear} starts afresh.
    \item \code{Shift+Tab} cycles the permission modes; \code{/permissions}
      edits allow and deny rules; \code{/sandbox} toggles the sandbox where
      supported \parencite{anthropic2025sandboxing}.
    \item \code{Esc} interrupts; \code{/rewind} rolls code and conversation
      back to a checkpoint; \code{/resume} reopens an earlier session.
    \item \code{claude --permission-mode plan} starts read-only, as Chapter 10
      does; \code{claude -p "prompt"} runs once and prints; \code{claude -c}
      continues the last session.
    \item \code{/memory} lists the rules files and opens the auto-memory
      folder: one Markdown note per file and an index loaded every session.
      Switch auto memory off there if notes should not stay on the machine
      (Chapter 7).
  \end{itemize}
  \source{code.claude.com/docs/en/commands, /interactive-mode (read 2 November 2026) and /memory (read 4 November 2026).}
\end{frame}

\begin{frame}{OpenCode: the controls}
  \begin{itemize}
    \item \code{/models} picks the model; the list comes from the providers
      in \code{opencode.json} and \code{/connect}.
    \item \code{Tab} switches between the Plan agent, which proposes without
      editing, and the Build agent, which edits (Chapter 10, Step 3).
    \item \code{Ctrl+T} cycles the model's variant, the provider's reasoning
      effort; \code{opencode run --variant high} sets it for one run.
    \item \code{Ctrl+X} then \code{S} opens the status view; \code{opencode
      stats} prints token use and cost per model and tool.
    \item \code{/compact} summarises; \code{/undo} and \code{/redo} step
      through the agent's changes; \code{Esc} interrupts.
    \item \code{opencode run "prompt"} runs without the interface;
      \code{-m ollama/gemma4:12b} picks a model, \code{--continue} the last
      session.
  \end{itemize}
  \source{opencode.ai/docs/cli and /keybinds; opencode --help, version 1.17.11.}
\end{frame}

\begin{frame}{Codex: the controls}
  \begin{itemize}
    \item \code{/model} chooses the model and its reasoning effort;
      \code{/status} shows the chat, how much context is used and the rate
      limits; \code{/compact} summarises.
    \item \code{/permissions} chooses between ``Ask for approval'',
      ``Approve for me'', ``Full access'' and the rules in
      \code{config.toml}. Start with ``Ask for approval''.
    \item From the command line: \code{-s read-only}, \code{workspace-write}
      or \code{danger-full-access} sets the sandbox; \code{-a on-request} or
      \code{never} sets when it asks.
    \item \code{/init} writes an \code{AGENTS.md} scaffold; \code{/review}
      reviews uncommitted changes.
    \item \code{codex exec "prompt"} runs once; \code{codex exec resume --last
      "next prompt"} continues it; \code{codex resume} reopens a session.
  \end{itemize}
  \source{developers.openai.com/codex (CLI pages, read 2 November 2026); codex --help, version 0.162.1.}
\end{frame}

\begin{frame}{Tokens, context and usage: what to watch}
  \begin{itemize}
    \item A token is about four characters, three quarters of a word. A
      200-line R script is one to two thousand tokens (Chapter 7).
    \item An agent re-sends the whole session at every step: the rules file,
      every file it read, every earlier turn. Prompts run to tens of
      thousands of tokens (Chapter 13).
    \item When the window fills, old context is cut and the agent forgets.
      Watch \code{/context}, \code{/status} or the status view; compact or
      start a new session at a natural break.
    \item Hosted plans meter usage in windows; the harness shows where you
      stand (\code{/usage}, \code{/status}). The figures of the day are in
      Appendix D.
    \item Local models have no meter. The cost is time: the wait before the
      first token grows with the prompt (Chapter 13).
  \end{itemize}
  \source{Chapters 7 and 13; Appendix D.}
\end{frame}

\begin{frame}[fragile]{Running an agent without a conversation}
  \begin{lstlisting}[style=shell]
claude -p "Describe data/raw/occurrences.csv; change nothing." > logs/describe.md
opencode run -m ollama/gemma4:12b "Describe data/raw/occurrences.csv" > logs/describe.md
codex exec "Describe data/raw/occurrences.csv" -o logs/describe.md
  \end{lstlisting}
  \begin{itemize}
    \item One prompt in, text out: the shape of a script step, and of the
      weekly review's assistant layer (Chapter 11).
    \item The permissions still apply. Keep the deny rules, and never use a
      ``skip permissions'' flag outside a container with no network: least
      privilege, the first of the Singapore Consensus principles for agents
      \parencite{casper2026singapore}.
    \item Save the output under \code{logs/} with the date, the model and the
      prompt: that is the provenance record of Chapter 9.
  \end{itemize}
  \source{The three tools' help, 2 November 2026; Chapters 9, 11 and 15.}
\end{frame}

\begin{frame}{Fencing the agent in: where each control sits}
  \fig[0.6\textheight]{ch15-fences}
  \source{Figure from Chapter 15 of the book. Permission prompts, the rules
  file, the harness configuration, the sandbox, a container, and a cut
  network: each has an honest limit.}
\end{frame}

\section{Before the session}

\begin{frame}[fragile]{A checklist that takes two minutes}
  \begin{lstlisting}[style=shell]
git --version && ollama --version && opencode --version   # or claude, codex
ollama list                              # the workshop's model is there
ollama ps                                # loaded, CONTEXT 32768 or more
cp -R workshop/session ~/survey && cd ~/survey
git init && git add -A && git commit -m "Before the agent session"
chmod -R a-w data/raw
opencode                                 # reads AGENTS.md and opencode.json
  \end{lstlisting}
  \begin{itemize}
    \item If the harness cannot see the model: check \code{opencode.json}
      names the tag \code{ollama list} shows, and that \code{tool\_call} is
      set.
    \item If the agent forgets the task: the context is too small on one side
      or the other.
    \item If a command is refused: that is the configuration working. Read
      the rule before widening it.
  \end{itemize}
  \source{Chapter 10, Steps 1 to 3; Chapter 8, the two settings.}
\end{frame}

\begin{frame}{Takeaways}
  \begin{itemize}
    \item Install, then verify: every tool prints its version, Ollama lists
      the model, the harness starts inside the project folder.
    \item Commit before the agent starts; read the diff after; keep a way
      back. Deny the Git commands that would remove it.
    \item Enlarge the context on both sides, runtime and harness; it is the
      setting that most often makes a local agent look broken.
    \item Every harness has the same five controls: model, effort, context
      and usage, permissions, and a one-shot mode. Learn them once.
    \item A refused command is the configuration doing its job. Widen a rule
      deliberately, not because the agent asked.
  \end{itemize}
\end{frame}

\begin{frame}[allowframebreaks]{References}
  \printbibliography[heading=none]
\end{frame}

\end{document}
